% GENERATED by tools/gen.py from the unit's vocabulary -- do not edit.
Ćwiczenia tego bloku są \emph{generowane} z listy słów tygodnia i sprawdzają pamięć słownictwa, nie gramatykę. Jeśli odpowiedź jest synonimem, a nie słowem z listy, sprawdź w tabeli powyżej.
\begin{cwiczeniadalej}
\exitem{Podaj rodzajnik: \dots{} \de{Zahl} (liczba)}{\artf{die} Zahl}
\exitem{Podaj rodzajnik: \dots{} \de{Alter} (wiek)}{\artn{das} Alter}
\exitem{Podaj rodzajnik: \dots{} \de{Jahr} (rok)}{\artn{das} Jahr}
\exitem{Podaj rodzajnik: \dots{} \de{Telefonnummer} (numer telefonu)}{\artf{die} Telefonnummer}
\exitem{Podaj rodzajnik: \dots{} \de{Adresse} (adres)}{\artf{die} Adresse}
\exitem{Podaj rodzajnik: \dots{} \de{Handy} (telefon komórkowy)}{\artn{das} Handy}
\exitem{Podaj rodzajnik: \dots{} \de{Formular} (formularz)}{\artn{das} Formular}
\exitem{Podaj rodzajnik: \dots{} \de{Punkt} (kropka; punkt)}{\artm{der} Punkt}
\exitem{Podaj rodzajnik: \dots{} \de{Bindestrich} (myślnik, łącznik)}{\artm{der} Bindestrich}
\exitem{Podaj rodzajnik: \dots{} \de{Unterstrich} (podkreślnik)}{\artm{der} Unterstrich}
\exitem{Podaj liczbę mnogą: \de{\artf{die} Zahl}}{\de{\artf{die} Zahlen}}
\exitem{Podaj liczbę mnogą: \de{\artf{die} Nummer}}{\de{\artf{die} Nummern}}
\exitem{Podaj liczbę mnogą: \de{\artf{die} Telefonnummer}}{\de{\artf{die} Telefonnummern}}
\exitem{Podaj liczbę mnogą: \de{\artf{die} Straße}}{\de{\artf{die} Straßen}}
\exitem{Podaj liczbę mnogą: \de{\artf{die} Postleitzahl}}{\de{\artf{die} Postleitzahlen}}
\exitem{Podaj liczbę mnogą: \de{\artn{das} Formular}}{\de{\artf{die} Formulare}}
\exitem{Podaj liczbę mnogą: \de{\artm{der} Punkt}}{\de{\artf{die} Punkte}}
\exitem{Podaj liczbę mnogą: \de{\artm{der} Bindestrich}}{\de{\artf{die} Bindestriche}}
\exitem{Jak to jest po niemiecku: jeden?}{\de{eins}}
\exitem{Jak to jest po niemiecku: dziewięć?}{\de{neun}}
\exitem{Jak to jest po niemiecku: osiemdziesiąt?}{\de{achtzig}}
\exitem{Jak to jest po niemiecku: dziewięćdziesiąt?}{\de{neunzig}}
\exitem{Jak to jest po niemiecku: liczba?}{\de{\artf{die} Zahl, Zahlen}}
\exitem{Jak to jest po niemiecku: numer?}{\de{\artf{die} Nummer, Nummern}}
\exitem{Jak to jest po niemiecku: adres?}{\de{\artf{die} Adresse, Adressen}}
\exitem{Jak to jest po niemiecku: kod pocztowy?}{\de{\artf{die} Postleitzahl, Postleitzahlen}}
\exitem{Jak to jest po niemiecku: literować?}{\de{buchstabieren}}
\exitem{Jak to jest po niemiecku: potrzebować?}{\de{brauchen}}
\exitem{Co znaczy \de{drei}?}{trzy}
\exitem{Co znaczy \de{neun}?}{dziewięć}
\exitem{Co znaczy \de{siebzig}?}{siedemdziesiąt}
\exitem{Co znaczy \de{hundert}?}{sto}
\exitem{Co znaczy \de{\artn{das} Alter}?}{wiek}
\exitem{Co znaczy \de{\artf{die} Nummer, Nummern}?}{numer}
\exitem{Co znaczy \de{\artf{die} Straße, Straßen}?}{ulica}
\exitem{Co znaczy \de{\artm{der} Bindestrich, Bindestriche}?}{myślnik, łącznik}
\exitem{Co znaczy \de{schreiben}?}{pisać}
\exitem{Co znaczy \de{fragen}?}{pytać}
\end{cwiczeniadalej}
